\documentclass[12pt,a4paper]{article}

\usepackage{fontspec}
\usepackage{polyglossia}
\setmainlanguage{english}
\setotherlanguage{russian}
\setmainfont{Times New Roman}
\newfontfamily\russianfont{Times New Roman}
\usepackage[a4paper,left=30mm,right=20mm,top=20mm,bottom=20mm]{geometry}
\usepackage{setspace}
\usepackage{microtype}
\usepackage{amsmath,amssymb}
\usepackage{booktabs}
\usepackage{longtable}
\usepackage{tabularx}
\usepackage{array}
\usepackage{enumitem}
\usepackage{caption}
\usepackage{graphicx}
\usepackage{xcolor}
\usepackage{fancyhdr}
\usepackage{hyperref}
\usepackage{xurl}

\definecolor{linkblue}{RGB}{25,76,130}
\hypersetup{
  colorlinks=true,
  linkcolor=black,
  citecolor=linkblue,
  urlcolor=linkblue,
  pdftitle={Experimental Detection of AI-Generated Images Using Spatial and Frequency-Domain Representations},
  pdfauthor={Anna Malyk Germanovna; Alisa Solovyeva Dmitrievna},
  pdfsubject={Course Project Report, HSE, 2026}
}

\setstretch{1.35}
\setlength{\parindent}{1.25cm}
\setlength{\parskip}{0pt}
\setlength{\emergencystretch}{2em}
\setlist{nosep,leftmargin=1.2cm}
\captionsetup{font=small,labelfont=bf}
\renewcommand{\arraystretch}{1.2}
\newcolumntype{Y}{>{\raggedright\arraybackslash}X}
\newcolumntype{R}{>{\raggedleft\arraybackslash}p{2.2cm}}

\pagestyle{fancy}
\fancyhf{}
\renewcommand{\headrulewidth}{0pt}
\cfoot{\thepage}

\begin{document}
\selectlanguage{english}
\renewcommand{\contentsname}{Contents}
\renewcommand{\refname}{References}

% The title-page identities follow CourseworkTemplateEng__2_.pdf supplied by the user.
\begin{titlepage}
\begin{center}
\textbf{NATIONAL RESEARCH UNIVERSITY}\\
\textbf{HIGHER SCHOOL OF ECONOMICS}

\vspace{1.0cm}
Faculty of Computer Science\\
Bachelor's Programme ``Applied Mathematics and Informatics''

\vfill
{\Large\textbf{Course Project Report}}\\[0.5cm]
{\large on the topic}\\[0.5cm]
{\Large\textbf{``Experimental Detection of AI-Generated Images\\
Using Spatial and Frequency-Domain Representations''}}

\vfill
\begin{tabularx}{\textwidth}{@{}Y Y@{}}
\textbf{Students:} & \textbf{Supervisor:}\\[0.3cm]
group \#БПАД234 & Makarov Artem Maksimovich\\
Anna Malyk Germanovna & Research Fellow,\\
Alisa Solovyeva Dmitrievna & Faculty of Computer Science, HSE\\
\end{tabularx}

\vfill
Moscow 2026
\end{center}
\end{titlepage}

\pagenumbering{arabic}
\setcounter{page}{2}
\tableofcontents
\clearpage

\section*{Annotation}
\addcontentsline{toc}{section}{Annotation}

This coursework studies a narrow and falsifiable version of AI-generated image detection:
binary classification of real COCO photographs and images produced by declared DANI generator
conditions. The main objective is not to maximise one attractive score, but to determine whether
an observed distinction survives dataset auditing, parent-disjoint splitting, matched training,
and common post-processing. The report therefore separates acquisition shortcuts, validation-only
model selection, locked internal-test evidence, robustness, and external transfer.

The selected high-resolution corpus contains 7,410 images from 1,482 parent groups. Every group
contains one real 1024×1024 COCO image and four generated counterparts: DALL-E~3 text-to-image,
SDXL text-to-image, SDXL image-to-image, and SDXL text-and-image-to-image. An original-byte audit
found no cross-label exact duplicate or cross-split integrity component, but it identified
label-dependent JPEG/PNG and colour-mode support. All selected files were therefore converted
losslessly to a canonical 1024×1024 RGB PNG corpus while retaining source and derived hashes. A
metadata-only validation control fell from ROC-AUC 0.713427 on original files to 0.537373 on the
canonical corpus, showing that canonicalisation removed most of this simple shortcut.

The primary experiment compared RGB and FFT-magnitude ResNet-50 models trained from random
initialisation at seeds 7, 17, and 42. Both representations received the same 384×384 common raster,
optimisation budget, parent-paired sampler, and Apple MPS device. FFT achieved mean validation
balanced accuracy 0.787556, compared with 0.657138 for RGB, and was selected by the predeclared
rule. Seed 17 was fixed in advance as the representative checkpoint. On the subsequently opened
1,110-image internal test, the frozen FFT model achieved ROC-AUC 0.893941, balanced accuracy
0.826014, and macro-F1 0.774670. The 95\% parent-group bootstrap interval for ROC-AUC was
[0.871820, 0.916324].

Robustness testing produced a strong negative result. With the checkpoint and validation threshold
unchanged, balanced accuracy fell to 0.518581 at JPEG quality 50, 0.476914 after a 0.75 resize,
0.478604 after a 0.50 resize, and 0.453829 under Gaussian blur. No post-test tuning was performed.
The results demonstrate a reproducible internal association in the declared DANI distribution,
but not a robust or general detector. A separately licensed Synthbuster + RAISE transfer corpus
was not materialised, so external performance remains not evaluated. The local web interface is
therefore kept behind a frozen-external-validation gate.

\vspace{0.6cm}
\begin{otherlanguage}{russian}
\section*{Аннотация}

В работе исследуется узкая и проверяемая постановка задачи распознавания изображений,
сгенерированных искусственным интеллектом: бинарная классификация реальных фотографий COCO и
изображений, полученных в заданных конфигурациях генераторов DANI. Основная цель заключается не в
получении одного максимального показателя, а в проверке того, сохраняется ли обнаруженная связь
после аудита датасета, разбиения по родительским группам, сопоставимого обучения и типовых
преобразований изображений.

Итоговый датасет содержит 7,410 изображений из 1,482 родительских групп. В каждой группе имеется
одна реальная фотография COCO и четыре синтетических варианта. Исходный аудит выявил зависимость
формата файла и цветового режима от класса, поэтому изображения были приведены без изменения
пиксельного содержания к единому формату RGB PNG 1024×1024 с сохранением хешей происхождения.
Модель, использующая только метаданные, снизила ROC-AUC с 0.713427 до 0.537373 после такой
канонизации.

Основной эксперимент сравнивал ResNet-50 для RGB и амплитудного спектра FFT при одинаковой
инициализации с нуля, одинаковом растре 384×384 и трёх случайных зернах. По среднему значению
сбалансированной точности на валидации было выбрано FFT-представление. На закрытой тестовой выборке
модель достигла ROC-AUC 0.893941 и сбалансированной точности 0.826014. При JPEG-сжатии, повторном
масштабировании и размытии качество резко снизилось. Поэтому результат подтверждает внутреннюю
связь в исследованном распределении, но не подтверждает устойчивость, переносимость на другие
источники или происхождение произвольного изображения.
\end{otherlanguage}
\clearpage

\section*{Keywords}
\addcontentsline{toc}{section}{Keywords}

AI-generated image detection; image forensics; frequency domain; fast Fourier transform; FFT
magnitude; dataset leakage; canonicalisation; group-disjoint split; ResNet-50; Apple MPS;
robustness; reproducible machine learning.

\section*{Abbreviations and notation}
\addcontentsline{toc}{section}{Abbreviations and notation}

\begin{tabularx}{\textwidth}{@{}p{3.2cm}Y@{}}
\toprule
Term & Meaning\\
\midrule
AI & Artificial intelligence.\\
FFT & Fast Fourier transform.\\
RGB & Red, green, and blue image channels.\\
ROC-AUC & Area under the receiver-operating-characteristic curve.\\
BAcc & Balanced accuracy: the arithmetic mean of real-class recall and AI-class recall.\\
FPR@TPR95 & False-positive rate at the first ROC operating point whose true-positive rate is at least 95\%.\\
ECE & Expected calibration error, calculated here with 15 confidence bins.\\
T2I / I2I / TI2I & Text-to-image, image-to-image, and text-and-image-to-image generation.\\
MPS & Apple's Metal Performance Shaders backend used by PyTorch.\\
Model score & An uncalibrated classifier output; not a probability or proof of origin.\\
\bottomrule
\end{tabularx}
\clearpage

\section{Introduction}

\subsection{Problem statement}

Modern text-to-image and image-to-image systems can produce visually convincing content at low
cost. This has created demand for tools that assist moderation, media verification, and research.
The broad phrase ``AI-image detector'', however, can conceal a much narrower statistical task. A
classifier observes correlations between labels and pixels in the available data. Those
correlations may arise from synthesis, but they may also arise from image dimensions, encoder,
file format, colour mode, prompt template, source website, or resizing history. A model that
separates one benchmark may therefore fail on an edited image or a different generator.

This distinction is important for academic validity. A high random-split score is not sufficient
when several synthetic images share the same real parent or semantic content. Likewise, a spectral
model may detect a resampling pipeline rather than an intrinsic property of synthesis. The present
work treats these failure modes as primary experimental questions. Attractive but confounded
results are retained as diagnostic evidence and are not relabelled as detector performance.

The target in this study is deliberately limited. It is binary classification between real COCO
photographs and four declared generated conditions selected from DANI. The work concerns still
images of general scenes, not face swaps, video deepfakes, authorship attribution, cryptographic
provenance, or factual verification. A positive decision means only ``AI-like under this frozen
experimental model and benchmark''. It does not establish how an arbitrary file was created.

\subsection{Aim, object, subject, and tasks}

The \textbf{aim} is to determine whether spatial RGB or frequency-domain FFT representation yields
stronger and reproducible internal discrimination under a common high-resolution raster, while
making dataset shortcuts, model-selection rules, uncertainty, and robustness explicit.

The \textbf{object of study} is a grouped collection of real and AI-generated images with declared
source relationships. The \textbf{subject of study} is the effect of image representation and
post-processing on the performance of binary classifiers trained under matched conditions.

The following tasks were defined:

\begin{enumerate}
  \item Review analogous RGB, frequency-domain, cross-generator, and robustness work.
  \item Audit a candidate corpus before downloading or training and establish licence and lineage.
  \item Materialise a high-resolution subset that fits the local storage boundary and preserves a
        demonstrable parent key.
  \item Detect exact duplicates, near-duplicate candidates, cross-split groups, and non-pixel class
        shortcuts.
  \item Construct a canonical 1024×1024 RGB PNG corpus and a parent-disjoint manifest for train,
        validation, and test.
  \item Compare metadata and radial FFT controls on validation.
  \item Train matched RGB and FFT ResNet-50 models from scratch for seeds 7, 17, and 42 using Apple
        MPS and identical budgets.
  \item Select a representation from validation aggregates only and use a representative seed
        declared before test inspection.
  \item Evaluate the frozen checkpoint once on clean and predeclared degraded test conditions,
        preserving group-aware uncertainty and generator-specific results.
  \item Provide a reproducible notebook and a constrained local interface without promoting the
        result to authentication or general provenance detection.
\end{enumerate}

\subsection{Research questions and hypotheses}

The primary question is: \emph{under a common 384×384 raster created from audited 1024×1024 source
images, does an FFT-magnitude ResNet-50 provide stronger validation discrimination than a matched
RGB ResNet-50, and does the selected signal survive the locked internal test and common
degradations?}

Three hypotheses were declared:

\begin{description}[style=nextline]
  \item[H1---representation.] FFT magnitude contains useful discriminative information after both
  classes have passed through the same canonical rasterisation, and its mean validation BAcc will
  be compared with the matched RGB result over three seeds.
  \item[H2---transfer.] A strong internal result does not by itself establish transfer to a new
  real-photo source or external generator families. This requires a separately locked
  Synthbuster + RAISE evaluation.
  \item[H3---robustness.] JPEG compression, resizing, and blur can materially change the learned
  signal even when the checkpoint and decision threshold remain fixed.
\end{description}

H1 and H3 can be evaluated locally. H2 remains unevaluated because the external licensed corpus is
not part of the materialised project. This missing result constrains the conclusion rather than
being treated as a zero metric or silently omitted.

\subsection{Contribution and report structure}

The contribution is methodological and empirical. The project combines lineage-aware data
selection, a canonical image-byte audit, parent-disjoint splitting, shortcut controls, a matched
multi-seed representation experiment, a hash-pinned terminal evaluator, and explicit negative
robustness evidence. It does not claim a new state of the art.

Section 2 reviews analogous approaches. Section 3 documents collection and auditing. Section 4
defines preprocessing, training, selection, and evaluation. Section 5 reports results and threats
to validity. Section 6 describes the research interface and its gate. Section 7 states the final
conclusion. The appendix records exact reproducibility identifiers and the use of AI assistance.

\section{Review of analogous works}

\subsection{Spatial neural classifiers}

Convolutional neural networks learn spatial filters directly from image pixels. Residual networks
introduced identity shortcuts that improve optimisation of deep architectures \cite{he2016resnet}.
ResNet-50 is a useful reference model here because it is widely documented, available in a stable
PyTorch implementation, and large enough to learn both semantic and textural signals. Its use does
not imply that residual networks are inherently forensic or that a deeper model guarantees better
transfer.

RGB detectors can exploit colour, texture, content, and rendering artefacts. The same flexibility
is a risk: the model can also learn the collection process. Prior synthetic-image detection work
has repeatedly shown that performance can change under a new generator or a new preprocessing
pipeline \cite{ricker2022diffusion,corvi2023intriguing}. Consequently, architecture choice must be
separated from evidence about dataset construction.

\subsection{Frequency-domain evidence}

For a common image raster $x$, the FFT representation in this work is based on shifted log
magnitude:

\begin{equation}
  F(x)=\log\left(1+\left|\operatorname{fftshift}(\operatorname{FFT2}(x))\right|\right).
\end{equation}

Phase is discarded. The magnitude emphasises the distribution of spatial frequencies, which can
reflect periodic textures, upsampling, denoising, or generator-specific reconstruction. Previous
work has reported spectral differences between real and synthetic imagery
\cite{corvi2023intriguing,corvi2023detection}. However, compression and resizing also modify the
spectrum. Therefore a successful FFT classifier may be detecting an acquisition or transformation
pipeline rather than a stable synthesis trace.

The comparison in this coursework applies RGB and FFT processing only after the same 384×384 common
raster is produced. The FFT tensor is normalised to a three-channel model input so that the same
ResNet-50 architecture can be used in both arms. This removes architecture size and direct source
geometry as explanations for a representation difference, while not eliminating all possible
source associations.

\subsection{Cross-generator and degradation benchmarks}

GenImage introduced explicit cross-generator and degraded-image tasks, including compression,
blur, and low resolution \cite{zhu2023genimage}. These tasks express two distinct questions. The
first asks whether a representation transfers beyond the generators used for training. The second
asks whether a signal survives an altered media pipeline. A detector may succeed at one and fail at
the other.

Synthbuster provides generated imagery from diffusion-model pipelines and was designed for
detector assessment \cite{bammey2024synthbuster}. A suitable transfer evaluation also requires a
separate real-photo source; RAISE is a raw-image forensic dataset that can serve this role under its
research conditions \cite{dangnguyen2015raise}. In this project these external sources are never
mixed into training or representation selection. Because they were not materialised locally, the
report makes no external numerical claim.

\subsection{Dataset bias and grouped evaluation}

The DANI dataset provides natural COCO images and synthetic variants across models, resolutions,
and generation protocols \cite{liu2025djudge,dani2026card}. Its parent-like structure is valuable
for paired research, but it creates a leakage risk if related images cross partitions. A random
image split can place a real parent in training and its generated counterpart in validation or
test. The apparent test then measures memorisation of content as well as synthesis association.

The solution used here is to assign a complete parent group to one split. Exact file hashes,
decoded-pixel hashes, and perceptual-hash candidates extend the boundary beyond a single metadata
field. Bootstrap intervals also resample complete parent groups instead of treating the five rows
in each group as independent observations.

\subsection{Gap addressed by the project}

Many educational projects begin with model training and inspect data problems only after metrics
look suspicious. This work reverses that order. The data source must pass provenance, licence,
lineage, geometry, duplicate, split, and shortcut gates before neural training. The model-selection
command cannot access internal test metrics. The terminal evaluator cannot train, change the
threshold, replace the checkpoint, or overwrite prior evidence. The resulting score is less broad
than a generic detector claim but more academically defensible.

\section{Collecting a dataset}

\subsection{Source selection and licence boundary}

DANI was chosen because it provides natural COCO images and multiple generated conditions with
explicit model, size, and generation-type metadata. The official card describes more than
445,000 images, several spatial resolutions, and generation modes including T2I, I2I, and TI2I
\cite{dani2026card}. It is distributed for non-commercial research under CC BY-NC 4.0. The local
workflow pins the dataset revision and records provenance instead of assuming that a moving online
dataset is immutable.

The complete dataset was not required for the declared comparison. A 40,GB materialisation cap was
used. Selection occurred from metadata before image bytes were read. The final subset deliberately
kept the exact 1024×1024 source level and four generator conditions for each recoverable COCO parent:
DALL-E~3 T2I, SDXL T2I, SDXL I2I, and SDXL TI2I. This produces a balanced parent structure even
though the binary class count is four synthetic rows for every real row.

\subsection{Lineage and parent key}

Metadata were first scanned without downloading image bytes. Candidate rows were accepted only
when their relationship to a common parent could be demonstrated by the DANI metadata workflow.
The parent COCO image identifier became the grouping key. Each retained group contains exactly one
real row and four synthetic rows. Groups that did not satisfy the required structure were excluded
rather than repaired by guessed filename similarity.

The final corpus contains 1,482 groups and 7,410 rows. Complete parent groups were allocated to
train, validation, and internal test. Table~\ref{tab:split} shows the result.

\begin{table}[ht]
\centering
\caption{Parent-disjoint DANI split.}
\label{tab:split}
\begin{tabular}{lrrrr}
\toprule
Split & Parent groups & Real & Synthetic & Total\\
\midrule
Train & 1,037 & 1,037 & 4,148 & 5,185\\
Validation & 223 & 223 & 892 & 1,115\\
Internal test & 222 & 222 & 888 & 1,110\\
\midrule
Total & 1,482 & 1,482 & 5,928 & 7,410\\
\bottomrule
\end{tabular}
\end{table}

No parent crosses a split. Within each split, all four generated conditions have the same count as
the real parent count. This makes generator-stratified real-versus-generator evaluation possible
without changing the real reference set.

\subsection{Original-byte integrity audit}

Downloaded files were decoded and checked against the materialisation manifest. The audit counted
rows, source cells, dimensions, formats, colour modes, exact file hashes, decoded-pixel hashes, and
perceptual-hash candidate links. It found:

\begin{itemize}
  \item zero exact duplicate groups in the selected rows;
  \item zero cross-label exact duplicate groups;
  \item zero cross-parent perceptual-hash links at Hamming distance at most four;
  \item zero integrity components crossing the proposed split;
  \item exact 1024×1024 geometry for every selected source image.
\end{itemize}

These findings support the chosen leakage boundary but do not prove the absence of every semantic
near-duplicate. Perceptual hashes are screening tools, and two visually related images can remain
outside a distance-four link. The known parent key is therefore still the primary independence
unit.

\subsection{Shortcut discovery and canonicalisation}

The same audit found that file container and decoded colour-mode support differed between labels.
Real and synthetic classes did not have identical JPEG/PNG and L/RGB support. A neural model could
learn codec residue or colour conversion rather than synthesis. Training on the original bytes was
therefore blocked even though duplicate and split checks passed.

Every selected source was decoded with orientation normalisation, converted to RGB, and written as
a 1024×1024 PNG without resizing. The derived manifest stores both source and output hashes. The
integrity audit was then repeated on the canonical corpus and verified uniform PNG/RGB support,
unchanged row/group structure, and no cross-split integrity component. The frozen training
manifest SHA-256 is:

\begin{center}
\small\path{3b88920920add51ef8c55b225817448b759c05a8d5bfc1cee11bc6befced2855}.
\end{center}

Canonicalisation does not make the corpus bias-free. PNG encoded size can still correlate with
image content or generator texture. It does, however, prevent format and colour mode from being a
trivial explicit separator.

\subsection{Metadata shortcut controls}

A class-balanced logistic regression used geometry, encoded size, container, and colour mode but
no decoded pixels. On the original validation files it achieved ROC-AUC 0.713427 and BAcc
0.660314. On the canonical files, ROC-AUC fell to 0.537373 and BAcc to 0.573430. Table
\ref{tab:metadata-control} summarises the result.

\begin{table}[ht]
\centering
\caption{Validation-only metadata shortcut controls.}
\label{tab:metadata-control}
\begin{tabular}{lrrr}
\toprule
Input & ROC-AUC & BAcc & Macro-F1\\
\midrule
Original files & 0.713427 & 0.660314 & 0.552612\\
Canonical RGB PNG & 0.537373 & 0.573430 & 0.509027\\
\bottomrule
\end{tabular}
\end{table}

The canonical value remains above 0.5, which is a warning that acquisition associations were
reduced rather than proven absent. The control is excluded from model selection because it does
not answer a pixel-representation question.

\section{Methodology, training, and evaluation protocol}

\subsection{Common 384×384 raster}

Every neural input begins from an audited exact 1024×1024 RGB PNG. The image is downsampled once to
384×384 with Lanczos interpolation. Upsampling is prohibited. Training may apply a horizontal flip
with probability 0.5 after the common raster is formed. Validation and test transformations are
deterministic.

The RGB arm normalises the three raster channels using the declared model transform. The FFT arm
converts the same common raster to grayscale, calculates shifted log magnitude, applies robust
normalisation, and replicates the result to three channels for ResNet-50. Thus any representation
difference occurs after shared image decoding and resizing.

\begin{center}
\fbox{\parbox{0.92\textwidth}{\centering
Audited 1024×1024 RGB PNG $\rightarrow$ one 384×384 Lanczos raster
$\rightarrow$ \{RGB normalisation \textbf{or} grayscale FFT log magnitude\}
$\rightarrow$ matched ResNet-50 $\rightarrow$ model score}}
\end{center}

\subsection{Model and optimisation contract}

Both primary arms use ResNet-50 with 23,512,130 trainable parameters and random initialisation.
The matched contract is:

\begin{itemize}
  \item seeds 7, 17, and 42 for both RGB and FFT;
  \item batch size 64;
  \item learning rate $10^{-4}$ and weight decay $10^{-4}$;
  \item maximum 15 epochs;
  \item early stopping patience four using validation loss;
  \item parent-paired binary sampler: one real and one uniformly selected synthetic counterpart
        from each sampled parent;
  \item Apple MPS device;
  \item no internal test scoring during training or selection.
\end{itemize}

The paired sampler prevents the four-to-one class ratio from turning each epoch into four times as
many synthetic observations. It also focuses optimisation on within-parent real--synthetic
differences while leaving the official evaluation distribution unchanged.

\subsection{Apple MPS verification and memory choice}

The local computer reported approximately 40.2~GB as the MPS recommended maximum allocation. A
direct synthetic ResNet-50 training-step benchmark at 384×384 and batch 16 measured 4.10 images/s
on CPU and 46.89 images/s on MPS on average: an observed 11.44× compute-loop speed-up. The benchmark
included host-to-device transfer and synchronised timing but excluded image decoding and data-loader
cost, so it is not an end-to-end training estimate.

Batch 64 on MPS sustained approximately 43 images/s and used about 20.37,GB of driver-allocated
memory in the recorded benchmark, below the recommended maximum and consistent with the requested
approximately 25,GB training budget. The study therefore uses 384×384 rather than the historical
128×128 raster. A larger raster would increase activation memory roughly with pixel area and reduce
batch size; 384×384 was retained as a measured compromise, not inferred only from nominal unified
memory.

\subsection{Baselines and practical RGB series}

Two auxiliary series were kept separate from the matched comparison. First, a deterministic
64-bin radial log-power FFT representation was fitted with class-balanced logistic regression. It
obtained validation ROC-AUC 0.862972, BAcc 0.803251, macro-F1 0.810515, FPR@TPR95 0.376682, and ECE
0.170572. This indicates pixel-frequency association but inadequate calibration and high
false-positive cost at high recall.

Second, an ImageNet-pretrained RGB ResNet-50 practical series ran over three seeds and produced mean
validation ROC-AUC 0.987991 and mean BAcc 0.951233. It demonstrates that the corpus and MPS path can
support a strong practical RGB signal. It is not used in the RGB-versus-FFT decision because FFT
has no matched ImageNet pretraining prior and the scientific comparison requires equal
initialisation.

\subsection{Validation-only selection}

The primary selection criterion is mean validation BAcc across all three seeds for each
representation. Exact ties are resolved by mean validation ROC-AUC and then lexicographic
representation name. Seed 17 is the representative checkpoint for whichever representation wins;
the best observed seed is never substituted. The validation threshold for that checkpoint is the
threshold that maximises validation BAcc.

The aggregation script refuses an incomplete representation/seed queue, unequal run signatures,
missing validation metrics, or internal test artifacts. The selection record explicitly states
that internal test metrics were not used. For the selected experiment, the threshold was
0.8226767778396606 and the checkpoint SHA-256 was:

\begin{center}
\small\path{0600e1e3346058ad826f61f37515022aad275b221cf9e7c14255d01b57be89c4}.
\end{center}

\subsection{Metrics and uncertainty}

The report includes ROC-AUC, BAcc, macro-F1, AI precision/recall, real recall, PR-AUC,
FPR@TPR95, ECE, Brier score, and confusion counts. Accuracy alone is secondary because the test
contains four synthetic images per real image.

At a fixed threshold $t$, balanced accuracy is

\begin{equation}
  \operatorname{BAcc}(t)=\frac{1}{2}\left(
  \frac{TP}{TP+FN}+\frac{TN}{TN+FP}\right).
\end{equation}

FPR@TPR95 is read from each condition's ROC curve and is descriptive; it does not replace the
validation-selected operating threshold. Clean uncertainty uses 2,000 percentile bootstrap draws.
Each draw resamples complete \texttt{group\_id} units, keeping the real row and all relevant
synthetic counterparts together.

\subsection{Terminal test and robustness gate}

The terminal evaluator validates the selection-record schema and status, the complete seed queue,
the manifest hash, checkpoint hash, representation, representative seed, preprocessing protocol,
and threshold. It refuses to proceed if the selected experiment already contains internal metrics
or if the output directory exists. The command exposes no training, checkpoint, representation, or
threshold options.

Seven conditions were declared before the terminal run: clean, JPEG Q95, JPEG Q75, JPEG Q50,
resize 0.75, resize 0.50, and Gaussian blur radius one. Every condition retains row-level manifest
provenance and the same frozen threshold. Results are not averaged across conditions.

\section{Results and discussion}

\subsection{Matched validation results}

Table~\ref{tab:seed-results} reports all six primary runs. There is no favourable-seed filtering.

\begin{table}[ht]
\centering
\caption{Validation metrics for matched from-scratch ResNet-50 runs.}
\label{tab:seed-results}
\small
\begin{tabular}{llrrrr}
\toprule
Representation & Seed & ROC-AUC & BAcc & Macro-F1 & FPR@TPR95\\
\midrule
FFT & 7  & 0.922203 & 0.843049 & 0.796153 & 0.336323\\
FFT & 17 & 0.897454 & 0.823430 & 0.772210 & 0.394619\\
FFT & 42 & 0.736341 & 0.696188 & 0.573436 & 0.843049\\
RGB & 7  & 0.732917 & 0.675448 & 0.601795 & 0.829596\\
RGB & 17 & 0.739031 & 0.692265 & 0.584206 & 0.843049\\
RGB & 42 & 0.633946 & 0.603700 & 0.472148 & 0.883408\\
\bottomrule
\end{tabular}
\end{table}

FFT outperformed RGB on every listed metric at the representation mean, but the FFT arm varied
substantially between seeds. Its mean ROC-AUC was 0.852000 with sample SD 0.100925; RGB achieved
0.701965 with SD 0.058985. Mean BAcc was 0.787556 for FFT and 0.657138 for RGB. Table
\ref{tab:aggregate-results} shows the aggregate.

\begin{table}[ht]
\centering
\caption{Three-seed validation aggregates (mean ± sample SD).}
\label{tab:aggregate-results}
\scriptsize
\begin{tabular}{lrrrr}
\toprule
Representation & ROC-AUC & BAcc & Macro-F1 & FPR@TPR95\\
\midrule
FFT & 0.852000 ± 0.100925 & 0.787556 ± 0.079732 & 0.713933 ± 0.122262 & 0.524664 ± 0.277266\\
RGB & 0.701965 ± 0.058985 & 0.657138 ± 0.047036 & 0.552717 ± 0.070326 & 0.852018 ± 0.028004\\
\bottomrule
\end{tabular}
\end{table}

H1 is supported within the validation protocol: FFT has the higher mean primary metric under equal
budget. The large variation prevents a stronger claim that FFT training is stable. The seed-17
checkpoint is representative by declaration, not because it is numerically best.

\subsection{Frozen clean internal test}

The clean test contains 222 real and 888 synthetic images from 222 parents. Table
\ref{tab:test-main} reports the frozen FFT seed-17 result and group-bootstrap intervals.

\begin{table}[ht]
\centering
\caption{Frozen clean internal-test result.}
\label{tab:test-main}
\begin{tabular}{lrr}
\toprule
Metric & Estimate & 95\% parent-group interval\\
\midrule
ROC-AUC & 0.893941 & [0.871820, 0.916324]\\
Balanced accuracy & 0.826014 & [0.800676, 0.851351]\\
Macro-F1 & 0.774670 & [0.752255, 0.798899]\\
FPR@TPR95 & 0.409910 & [0.337838, 0.472973]\\
\bottomrule
\end{tabular}
\end{table}

At threshold 0.822677, the confusion counts were TN=181, FP=41, FN=145, and TP=743. Real recall
was 0.815315 and AI recall 0.836712. ECE was 0.030147 and the Brier score was 0.089432. The clean
test ROC-AUC is close to the representative checkpoint's validation ROC-AUC of 0.897454, so there
is no large aggregate validation-to-test collapse. This agreement is internal to the same source
and preprocessing distribution.

\subsection{Generator-condition heterogeneity}

Each generator row in Table~\ref{tab:generator} compares all 222 real images with 222 generated
images from one condition.

\begin{table}[ht]
\centering
\caption{Clean internal test by generated condition.}
\label{tab:generator}
\small
\begin{tabular}{lrrrr}
\toprule
Condition & ROC-AUC & BAcc & Macro-F1 & FPR@TPR95\\
\midrule
DALL-E 3 T2I & 0.947772 & 0.889640 & 0.889027 & 0.157658\\
SDXL T2I & 0.962747 & 0.896396 & 0.895711 & 0.108108\\
SDXL I2I & 0.832299 & 0.754505 & 0.753593 & 0.481982\\
SDXL TI2I & 0.832948 & 0.763514 & 0.762877 & 0.509009\\
\bottomrule
\end{tabular}
\end{table}

The model is much stronger on the T2I conditions than on I2I/TI2I. This difference matters because
an aggregate score is dominated by four synthetic sources and can obscure the weakest condition.
The finding is consistent with the possibility that image-conditioned generation preserves more
natural-image frequency structure. It is an interpretation, not a causal demonstration.

\subsection{Robustness results}

Table~\ref{tab:robustness} shows the unchanged checkpoint and threshold under all predeclared
conditions.

\begin{table}[ht]
\centering
\caption{Aggregate robustness results.}
\label{tab:robustness}
\small
\begin{tabular}{lrrrrr}
\toprule
Condition & ROC-AUC & BAcc & Macro-F1 & Real recall & AI recall\\
\midrule
Clean & 0.893941 & 0.826014 & 0.774670 & 0.815315 & 0.836712\\
JPEG Q95 & 0.822234 & 0.716779 & 0.736664 & 0.500000 & 0.933559\\
JPEG Q75 & 0.723521 & 0.559685 & 0.561415 & 0.144144 & 0.975225\\
JPEG Q50 & 0.698736 & 0.518581 & 0.484081 & 0.040541 & 0.996622\\
Resize 0.75 & 0.380828 & 0.476914 & 0.447567 & 0.027027 & 0.926802\\
Resize 0.50 & 0.416144 & 0.478604 & 0.262398 & 0.833333 & 0.123874\\
Gaussian blur & 0.430051 & 0.453829 & 0.362177 & 0.585586 & 0.322072\\
\bottomrule
\end{tabular}
\end{table}

Even high-quality JPEG Q95 reduced BAcc by 0.109235. At Q75 and Q50, the fixed operating point
shifted strongly toward the AI class: real recall fell to 0.144144 and 0.040541. The model therefore
flags almost every image as synthetic after stronger JPEG compression. Resize 0.75 produced real
recall 0.027027, whereas resize 0.50 reversed behaviour and reduced AI recall to 0.123874. Blur
reduced ROC-AUC to 0.430051 and BAcc below chance.

These are not small calibration shifts. Several transformations change score ordering itself, as
shown by ROC-AUC below 0.5. H3 is supported. The checkpoint is brittle to the media operations that
an uploaded web image commonly undergoes. This negative result is central to the conclusion and is
not used to justify post-test augmentation or a replacement threshold.

\subsection{Comparison with controls}

The canonical metadata control had ROC-AUC 0.537373, the radial FFT logistic baseline 0.862972 on
validation, the matched FFT ResNet-50 mean 0.852000 across validation seeds, and the frozen neural
checkpoint 0.893941 on internal test. These values answer different questions. The control measures
non-pixel acquisition association. The radial baseline measures a fixed 64-bin spectral summary.
The multi-seed neural series supports representation selection. The terminal test measures one
predeclared checkpoint after selection. They must not be ranked as if they were exchangeable runs
on one common test set.

The pretrained RGB practical series achieved a stronger validation aggregate than the matched
from-scratch series. It is excluded from H1 because its initialisation contains an unmatched prior.
The result nevertheless suggests that representation is not the only determinant of performance;
pretraining and learned semantic features may dominate within the same corpus.

\subsection{Threats to validity}

\textbf{Internal validity.} Parent-disjoint splitting, hash audits, and a common raster reduce known
leakage and geometry shortcuts. They cannot exclude every semantic near-duplicate or correlation
between content category and generator condition. The remaining canonical metadata control is
weakly above chance.

\textbf{Statistical validity.} Three seeds provide limited evidence about neural variability. FFT
seed 42 was substantially weaker, so the mean and SD are both reported. Bootstrap intervals account
for parent dependence but do not account for uncertainty in dataset construction or generator
selection.

\textbf{Construct validity.} The benchmark label is not ground-truth provenance for arbitrary
uploads. The model distinguishes declared collections. It may use generator, codec, content, or
transformation-specific cues that are legitimate predictors inside DANI but irrelevant elsewhere.

\textbf{External validity.} No cross-dataset evaluation was completed. DANI has only one real-photo
source in the selected subset and four generated conditions. The result cannot be generalised to
social-media screenshots, scanned documents, edited photographs, new generators, or unknown
acquisition pipelines.

\textbf{Operational validity.} Robustness failures directly contradict deployment. A public service
would also require security review, upload limits, privacy policy, monitoring, and calibrated
decision costs. Those engineering controls cannot compensate for missing model evidence.

\section{Prototype inference interface}

\subsection{Purpose and current state}

The repository contains a small local web shell for research demonstrations. It is intentionally
not a public authenticity service. The interface can decode an uploaded image, apply a frozen
experiment transform, and return a model score together with the threshold, representation,
experiment name, and checkpoint hash.

The application refuses to serve an arbitrary experiment directory. It requires a separate model
selection record with schema \path{ai_image_detector_model_selection_v1} and status
\path{frozen_external_validated}. The current DANI selection record has status
\path{validation_selected_pending_external_validation}. Therefore the internally evaluated
checkpoint is not unlocked for the user-facing classifier.

\subsection{Safety and wording}

If a future externally validated record is created, the interface must:

\begin{enumerate}
  \item verify the checkpoint SHA-256 and preprocessing contract before inference;
  \item reject malformed, decompression-bomb, or excessively large uploads;
  \item avoid retaining uploaded images;
  \item display ``model score'', not ``probability of AI generation'';
  \item use ``AI-like under this experimental model'' and ``not AI-like under this experimental
        model'' rather than ``fake'', ``authentic'', or ``proven'';
  \item display tested generators, external-validation status, and known transformation failures.
\end{enumerate}

The gate is a deliberate project result. Enabling the best internal model despite failed
robustness would make the website more visually complete but academically misleading. The current
shell demonstrates the software boundary without converting a limited experiment into a consumer
claim.

\subsection{Reproducible software layout}

The repository keeps two notebooks. The first reads the DANI integrity audit, controls, six matched
validation runs, aggregate, and frozen selection record. The second reads the immutable terminal
evaluation and reproduces clean, generator-specific, and robustness tables. Both notebooks execute
from saved machine-readable artifacts and contain assertions for the main protocol gates.

The strict evaluator is covered by tests for test-informed selection rejection, incomplete seed
queue rejection, checkpoint mismatch, manifest mismatch, provenance-preserving predictions,
group-level intervals, and overwrite refusal. At the final snapshot, the full suite contained 208
passing tests.

\section{Conclusion}

This coursework investigated whether frequency-domain representation improves controlled
real-versus-AI image classification. The work began from a general detector idea but adopted a
strictly narrower claim after finding dataset shortcuts in earlier material. The final DANI study
was therefore gated by licence, lineage, exact geometry, duplicate checks, parent-disjoint
splitting, and acquisition-metadata controls before neural training.

The canonical corpus contains 7,410 1024×1024 RGB PNG images from 1,482 parents. Canonicalisation
reduced metadata-only validation ROC-AUC from 0.713427 to 0.537373. This is evidence that simple
file-support association was materially reduced, although not proven absent.

The matched from-scratch comparison completed for RGB and FFT across seeds 7, 17, and 42. FFT had
the higher mean validation BAcc, 0.787556 versus 0.657138, and was selected according to the frozen
rule. The predeclared seed-17 checkpoint achieved clean internal-test ROC-AUC 0.893941, BAcc
0.826014, and macro-F1 0.774670. Parent-group bootstrap intervals excluded chance for the principal
metrics. H1 is therefore supported within the declared DANI internal distribution.

The same model failed the robustness requirement. JPEG compression caused the fixed threshold to
classify most real images as synthetic; resizing and blur changed score ordering and reduced BAcc
to approximately chance or below. This supports H3 and falsifies any claim that the checkpoint is
ready for uploaded web images. No threshold, architecture, augmentation, or checkpoint was changed
after observing these results.

H2 remains unresolved. A separately licensed Synthbuster + RAISE corpus was not materialised, so
the project has no cross-dataset transfer number. The correct final statement is limited: an FFT
ResNet-50 learned a reproducible internal association in an audited DANI subset, but the signal is
generator-dependent, brittle to common transformations, and not externally validated. The model
score is not a calibrated probability or proof of provenance, and the web inference gate remains
closed.

The academic value lies in the full result, including negative evidence. The project produced a
reproducible dataset contract, matched multi-seed training, a frozen test, uncertainty intervals,
and a documented failure under realistic degradation. This is preferable to presenting a strong
validation score as a universal detector.

\clearpage
\begin{thebibliography}{99}

\bibitem{he2016resnet}
K. He, X. Zhang, S. Ren, and J. Sun,
``Deep Residual Learning for Image Recognition,''
\emph{Proceedings of the IEEE Conference on Computer Vision and Pattern Recognition},
pp. 770--778, 2016.
\href{https://openaccess.thecvf.com/content_cvpr_2016/html/He_Deep_Residual_Learning_CVPR_2016_paper.html}{Available online}.

\bibitem{ricker2022diffusion}
J. Ricker, S. Damm, T. Holz, and A. Fischer,
``Towards the Detection of Diffusion Model Deepfakes,'' arXiv:2210.14571, 2022.
\href{https://arxiv.org/abs/2210.14571}{Available online}.

\bibitem{corvi2023intriguing}
R. Corvi, D. Cozzolino, G. Poggi, K. Nagano, and L. Verdoliva,
``Intriguing Properties of Synthetic Images: From Generative Adversarial Networks to Diffusion
Models,'' arXiv:2304.06408, 2023.
\href{https://arxiv.org/abs/2304.06408}{Available online}.

\bibitem{corvi2023detection}
R. Corvi, D. Cozzolino, G. Zingarini, G. Poggi, K. Nagano, and L. Verdoliva,
``On the Detection of Synthetic Images Generated by Diffusion Models,''
\emph{IEEE International Conference on Acoustics, Speech and Signal Processing}, 2023.
\href{https://arxiv.org/abs/2211.00680}{Available online}.

\bibitem{zhu2023genimage}
M. Zhu, H. Chen, Q. Yan, X. Huang, G. Lin, W. Li, Z. Tu, H. Hu, J. Hu, and Y. Wang,
``GenImage: A Million-Scale Benchmark for Detecting AI-Generated Image,''
\emph{Advances in Neural Information Processing Systems}, vol. 36, 2023.
\href{https://proceedings.neurips.cc/paper_files/paper/2023/hash/f4d4a021f9051a6c18183b059117e8b5-Abstract-Datasets_and_Benchmarks.html}{Available online}.

\bibitem{bammey2024synthbuster}
Q. Bammey, ``Synthbuster: Towards Detection of Diffusion Model Generated Images,''
\emph{IEEE Open Journal of Signal Processing}, vol. 5, 2024.
doi: \href{https://doi.org/10.1109/OJSP.2023.3337714}{10.1109/OJSP.2023.3337714}.

\bibitem{dangnguyen2015raise}
D.-T. Dang-Nguyen, C. Pasquini, V. Conotter, and G. Boato,
``RAISE: A Raw Images Dataset for Digital Image Forensics,''
\emph{Proceedings of the 6th ACM Multimedia Systems Conference}, pp. 219--224, 2015.
doi: \href{https://doi.org/10.1145/2713168.2713194}{10.1145/2713168.2713194}.

\bibitem{liu2025djudge}
R. Liu, Z. Lyu, W. Zhou, and S.-K. Ng,
``D-Judge: How Far Are We? Assessing the Discrepancies Between AI-synthesized Images and Natural
Images through Multimodal Guidance,'' \emph{ACM International Conference on Multimedia}, 2025.
\href{https://arxiv.org/abs/2412.17632}{Available online}.

\bibitem{dani2026card}
Renyang Liu et al., ``DANI: Discrepancy Assessing for Natural and AI Images,'' dataset card,
revision pinned by the local provenance record, accessed 2026.
\href{https://huggingface.co/datasets/Renyang/DANI}{Available online}.

\bibitem{pytorch}
A. Paszke et al., ``PyTorch: An Imperative Style, High-Performance Deep Learning Library,''
\emph{Advances in Neural Information Processing Systems}, vol. 32, 2019.
\href{https://papers.neurips.cc/paper/9015-pytorch-an-imperative-style-high-performance-deep-learning-library}{Available online}.

\end{thebibliography}

\clearpage
\appendix
\section{Reproducibility record}

\begin{longtable}{@{}p{4.7cm}p{10.2cm}@{}}
\toprule
Item & Recorded value\\
\midrule
\endfirsthead
\toprule
Item & Recorded value\\
\midrule
\endhead
Canonical manifest & \path{artifacts/audits/dani_rgb1024_integrity_v1/training_manifest.csv}\\
Manifest SHA-256 & \path{3b88920920add51ef8c55b225817448b759c05a8d5bfc1cee11bc6befced2855}\\
Selection record & \path{artifacts/aggregates/dani_rgb_fft_scratch_validation_v1/prototype_selection_record.json}\\
Selected experiment & \path{artifacts/experiments/dani_fft_scratch_seed17_validation_v1}\\
Checkpoint SHA-256 & \path{0600e1e3346058ad826f61f37515022aad275b221cf9e7c14255d01b57be89c4}\\
Representation & FFT magnitude, representative seed 17\\
Raster & \path{highres_square_crop_384_v1}; one 384×384 Lanczos downsample from audited 1024×1024 RGB PNG\\
Model & ResNet-50 from random initialisation; 23,512,130 trainable parameters\\
Optimisation & Batch 64; learning rate $10^{-4}$; weight decay $10^{-4}$; maximum 15 epochs; patience four\\
Device & Apple MPS; PyTorch 2.13.0 recorded at launch\\
Threshold & 0.8226767778396606, selected on validation BAcc\\
Internal test & 1,110 images from 222 parents; 222 real and 888 synthetic\\
Uncertainty & 2,000 percentile bootstrap resamples; unit = \path{group_id}; seed 20260829\\
Terminal output & \path{artifacts/evaluations/dani_fft_seed17_internal_test_v1}\\
Terminal status & \path{completed_no_post_test_tuning}\\
Notebooks & \path{notebooks/01_internal_training.ipynb}; \path{notebooks/02_external_validation_and_results.ipynb}\\
\bottomrule
\end{longtable}

\section{Interpretation checklist}

Before a future result is used in the interface, all answers must be affirmative:

\begin{enumerate}
  \item Was the model selected without inspecting the relevant test or external labels?
  \item Did all compared representations receive the same raster and optimisation budget?
  \item Were parent, exact-hash, decoded-pixel-hash, and declared near-duplicate boundaries audited?
  \item Was the decision threshold selected on validation and then frozen?
  \item Were correlated rows resampled as complete groups?
  \item Were all declared seeds reported rather than only the strongest seed?
  \item Were generator-specific and degradation results kept separate from the clean aggregate?
  \item Was external transfer evaluated on a separately audited and licensed corpus?
  \item Does the interface describe a model score and its limitations rather than provenance?
\end{enumerate}

The current project answers ``no'' to item 8. This is why the website remains gated despite a
completed internal result.

\section{Disclosure of AI assistance}

AI assistance was used to help inspect repository artifacts, draft and refactor code, execute
tests, organise the report, and check consistency between machine-readable metrics and prose. The
dataset was not labelled by a language model, and reported numerical results were read from local
experiment artifacts rather than generated from prose. The evaluator, notebooks, report, and
citations were reviewed against the repository and supplied templates. Responsibility for the
experimental design, final claims, source use, and submitted text remains with the authors.

\end{document}
